\section{Related Work}

\noindent\textbf{Retrospective agent memory benchmarks.}
A large body of work evaluates whether agents can retrieve or recall information from long dialogues and multi-session histories. LoCoMo~\cite{locomo2024} stresses long-term conversational memory with questions grounded in prior turns; LongMemEval~\cite{longmemeval2024} and related suites probe retention, updates, and temporal reasoning over extended agent logs. Method papers respond with hierarchical summaries, graph memories, write--read controllers, and RAG variants~\cite{amem2025,memorybank2023,memgpt2023}. These benchmarks and systems optimize $\mathrm{sim}(e,q)$---whether the right past fact surfaces under a question \emph{now}---not $\mathrm{sat}(\varphi,e_t)$, i.e., whether a deferred trigger is satisfied at step $t$. They do not require executing a previously announced action at a future cue while concurrent tasks continue.

\noindent\textbf{Prospective memory in cognitive science.}
Prospective memory research studies how humans form delayed intentions, monitor for cues, and execute actions under ongoing-task load~\cite{einstein2005prospective,smith2003demand}. The Virtual Week paradigm~\cite{rendell2000virtual} embeds irregular time- and event-based tasks in a simulated seven-day schedule and scores whether participants act when due---a design that PM-Bench~\cite{pmbench2026} adapts for LLM agents. We inherit PM-Bench's due-set framing and diagnostic slices (channel hit, cross-day, update-sensitive) and treat them as the primary yardstick for conditional memory.

\noindent\textbf{Prospective memory benchmarks and scaffolds.}
PM-Bench~\cite{pmbench2026} is the first peer-reviewed benchmark to operationalize LLM prospective memory with Set-F1 over due sets $D_t$, hidden channels, lures, and lifecycle updates. It evaluates eight scaffolds: Single (no external memory), Todo-ledger (in-context pending-intent list), Optional heartbeat (model-initiated periodic self-checks), Auto-heartbeat at 60m/30m intervals, and Hierarchical union-query (coordinator + sub-agents that union channel reads before acting). Reported findings include: optional heartbeat leads macro Set-F1 ($65.1\%$) but still under-hits channel-required tasks; auto-heartbeat improves update-sensitive recall at the cost of FP inflation; hierarchical union-query maximizes queries yet converts poorly to correct $\hat{D}_t$; todo ledgers approximate dual storage but lack typed lifecycle APIs consumed by monitor and scorer modules. ProMem instantiates conditional memory \emph{on} this benchmark: PIS replaces unstructured ledger text with typed $(\varphi,\alpha,\sigma)$ records; $\pi_{\mathrm{mon}}$ replaces undifferentiated periodic wakes with intention-conditioned channel polling; $\pi_{\mathrm{due}}$ adds lifecycle and lure penalties absent from heartbeat-only scaffolds.

\noindent\textbf{Related memory lines (retrospective / hierarchical).}
Reflective memory management for personalized dialogue~\cite{inprospect2025} shares ``prospect/retrospect'' vocabulary but targets conversational QA with reflective write/read---not Set-F1 over hidden-channel due sets. HiAgent~\cite{hiagent2025} optimizes hierarchical \emph{working-memory} context control for long-horizon tool tasks rather than PM-Bench-style when-to-act scoring. Prompt-chaining recall frameworks~\cite{promptchain2025} implement reminder engineering without PIS lifecycles or lure-aware due-set scoring. These lines complement but do not substitute typed conditional memory on PM-Bench.

\noindent\textbf{Industrial periodic monitoring.}
Production agent stacks (OpenClaw, heartbeat-agent-framework, and similar cron-style loops, 2025--2026) overlap conceptually with PM-Bench's auto-heartbeat baselines: periodic environment wakes without intention-conditioned query budgets. We treat them as engineering instantiations of the same design axis PM-Bench already measured; ProMem's claim is \emph{intent-conditioned} monitoring and due-set precision control relative to fixed-interval polling, to be validated empirically rather than assumed.

\noindent\textbf{Positioning.}
Our work is complementary to retrospective memory research: episodic stores remain useful for answering history questions, but they are insufficient as the sole substrate for deferred execution. Conditional memory and ProMem target the orthogonal axis isolated and benchmarked by PM-Bench---triggering, updating, and suppressing intentions over time---by instantiating a typed scaffold rather than rediscovering the problem.
